% ============================================================
%  MODÈLE DE CHAPITRE — À COPIER
%
%  Utilisation :
%    1. Copiez ce fichier vers chapitres/ch3.tex
%    2. Remplacez le numéro et le titre du chapitre.
%    3. Remplacez les numéros de sections (3.1, 3.1.1, ...).
%    4. Ajoutez \input{chapitres/ch3} dans main.tex.
%    5. Compilez deux fois avec LuaLaTeX.
%
%  Ce fichier n'est pas inclus dans main.tex.
% ============================================================

\chapitre{3}{Titre du chapitre 3}

\begin{objectifs}
\begin{itemize}
  \item [Objectif 1 : à compléter]
  \item [Objectif 2 : à compléter]
  \item [Objectif 3 : à compléter]
\end{itemize}
\end{objectifs}

\soussection{3.1}{Titre de la section 3.1}

Rédigez ici l'introduction de la section.

\soussoussection{3.1.1}{Titre de la sous-section 3.1.1}

Rédigez ici le premier point détaillé.

\soussoussection{3.1.2}{Titre de la sous-section 3.1.2}

Rédigez ici le deuxième point détaillé.

\soussection{3.2}{Titre de la section 3.2}

Rédigez ici la suite du chapitre.

\begin{resume}
Résumez en quelques lignes les points essentiels du chapitre 3.
\end{resume}
